\begin{exercise}{9.1}{4}
    Sei $E$ eine reelle affine Ebene. Man finde alle Kollineationen $\mathbb V E \tilde{\to}
    \mathbb V E$ die die unendlich ferne Gerade punktweise festhalten. Man begründe, warum
    es keine anderen geben kann.
\end{exercise}

Jede solche Kollineation bildet insbesondere
\(\mathbb P\vec E\) bijektiv auf sich selbst ab und schränkt daher zu
einer Kollineation \(\varphi:E\to E\)ein. Nach LA1 3.3.1 sind dies genau die bijektiven affinen Abbildungen,
also die Affinitäten von \(E\).

Nach Wahl eines Ursprungs können wir eine solche Affinität in der Form
\[
    \varphi(x)=Lx+b \quad \text{mit} \quad L\in\operatorname{GL}(\vec E), \quad b\in\vec E
\]
schreiben.

Sei nun \(g=p+\mathbb R\vec v\) eine affine Gerade. Ihre projektive Vervollständigung ist \(g\sqcup\{[\vec v]\}\).
Da die unendlich ferne Gerade punktweise festgehalten wird, muss der
unendlich ferne Punkt \([\vec v]\) unter der Kollineation fest bleiben.
Daher hat die Bildgerade \(\varphi(g)\) dieselbe Richtung wie \(g\).

Da
\[
    \varphi(p+t\vec v)
    =
    Lp+tL\vec v +b,
\]
hat \(\varphi(g)\) den Richtungsraum \(\mathbb R L\vec v\).
Somit muss für jedes \(\vec v\neq0\) gelten
\[
    \mathbb R L\vec v=\mathbb R\vec v.
\]
Der lineare Anteil \(L\) stabilisiert also jeden eindimensionalen
Untervektorraum von \(\vec E\).

Wir zeigen, dass daraus \(L=\lambda\operatorname{id}_{\vec E} \)
für ein \(\lambda\in\mathbb R^\times\) folgt. Wähle eine Basis
\[
    e_1=(1,0),
    \qquad
    e_2=(0,1).
\]
Da die Räume \(\mathbb Re_1\) und \(\mathbb Re_2\)
stabilisiert werden, gibt es \(\lambda,\mu\in\mathbb R^\times\) mit
\[
    Le_1=\lambda e_1,
    \qquad
    Le_2=\mu e_2.
\]
Da auch \(\mathbb R(e_1+e_2)\)
stabilisiert wird, muss \(L(e_1+e_2)=\lambda e_1+\mu e_2\)
ein Vielfaches von \(e_1+e_2\) sein. Daher gilt
\[
    \lambda=\mu \quad \Longrightarrow \quad L=\lambda\operatorname{id}_{\vec E}.
\]
mit \(\lambda\neq0\).

Umgekehrt hat jede Affinität der Form
\[
    \varphi(x)=\lambda x+b,
    \qquad
    \lambda\in\mathbb R^\times,
    \quad
    b\in\vec E,
\]
die gewünschte Eigenschaft, denn \(\mathbb R(\lambda\vec v)=\mathbb R\vec v\)
für jeden \(\vec v\neq0\). Sie hält also jeden unendlich fernen Punkt
fest.

Damit sind die gesuchten Kollineationen genau die Affinitäten
\[
    \boxed{
        \varphi(x)=\lambda x+b,
        \qquad
        \lambda\in\mathbb R^\times,
        \quad
        b\in\vec E.
    }
\]